%%%%%%%%%%%%%%%%%%%%%%%%%%%%%%%%%%%%%%%%%%%%%%%%%%%%
% LECTURE 21
% 
% LECTURE: THERMAL MANAGEMENT IN FORMULA ONE
%%%%%%%%%%%%%%%%%%%%%%%%%%%%%%%%%%%%%%%%%%%%%%%%%%%%
\section*{Thermal Management in Formula One}
\addcontentsline{toc}{section}{Thermal Management in Formula One}

%%%%%%%%%%%%%%%%%%%%%%%%%%%%%%%%%%%%%%%%%%%%%%%%%%%%
\subsection*{Introduction to Thermal Management in Racing}
\addcontentsline{toc}{subsection}{Introduction to Thermal Management in Racing}
%%%%%%%%%%%%%%%%%%%%%%%%%%%%%%%%%%%%%%%%%%%%%%%%%%%%

Thermal management is a critical aspect of Formula One car design and performance. This lecture covers how thermal engineers at Red Bull Racing approach thermal challenges, solve complex problems under strict time constraints, and apply fundamental heat transfer principles in high-performance environments.

%%%%%%%%%%%%%%%%%%%%%%%%%%%%%%%%%%%%%%%%%%%%%%%%%%%%
\subsection*{Role of Thermal Engineers in Formula One}
\addcontentsline{toc}{subsection}{Role of Thermal Engineers in Formula One}
%%%%%%%%%%%%%%%%%%%%%%%%%%%%%%%%%%%%%%%%%%%%%%%%%%%%

The principal thermal systems engineer at Red Bull Racing is responsible for thermal management of the entire car, which is built from scratch every year. Key responsibilities include:

\begin{itemize}
    \item Managing thermal aspects of all car components
    \item Handling 50-100 different thermal tasks annually
    \item Solving problems within tight timeframes (same race weekend to several months)
    \item Designing cooling solutions for critical systems:
    \begin{itemize}
        \item Braking systems
        \item Radiator systems
        \item Gearboxes
        \item Tire management
        \item Sensors and electronics
    \end{itemize}
\end{itemize}

Formula One braking systems generate extreme thermal loads:
\begin{itemize}
    \item Peak power of approximately 1 megawatt during braking events
    \item Requires advanced cooling solutions to manage temperature
\end{itemize}

%%%%%%%%%%%%%%%%%%%%%%%%%%%%%%%%%%%%%%%%%%%%%%%%%%%%
\subsection*{Case Study: Pressure Sensor Thermal Management}
\addcontentsline{toc}{subsection}{Case Study: Pressure Sensor Thermal Management}
%%%%%%%%%%%%%%%%%%%%%%%%%%%%%%%%%%%%%%%%%%%%%%%%%%%%

\subsubsection*{Problem Definition}
\addcontentsline{toc}{subsubsection}{Problem Definition}

A pressure sensor on the car exceeded its thermal operating requirements:
\begin{itemize}
    \item Sensor temperature must remain below 5°C above ambient temperature
    \item Current sensor consistently operates above this requirement
    \item Incorrect temperature leads to inaccurate pressure readings
    \item Affects understanding of car performance (e.g., downforce measurements)
\end{itemize}

\subsubsection*{Project Requirements}
\addcontentsline{toc}{subsubsection}{Project Requirements}

The thermal management solution needed to:
\begin{itemize}
    \item Match track data via simulation
    \item Simulate multiple potential solutions
    \item Be lightweight and financially viable
    \item Use readily available materials for immediate implementation
    \item If immediate solutions weren't viable, propose R\&D for new cooling techniques or materials
\end{itemize}

\subsubsection*{Sensor Configuration}
\addcontentsline{toc}{subsubsection}{Sensor Configuration}

The pressure sensor assembly included:
\begin{itemize}
    \item Hot air flow heating the sensor from above
    \item Sensor mounted in a casing
    \item Rapid prototyping leads
    \item Foam insulation
    \item Anti-vibration mounting system
    \item Rubber sock to prevent vibrational damage
    \item Carbon fiber composite floor structure:
    \begin{itemize}
        \item Upper carbon fiber layer
        \item Core sandwich structure (polymer-based Rohacell or aluminum honeycomb)
        \item Lower carbon fiber layer
    \end{itemize}
    \item Airflow under the floor structure
\end{itemize}

%%%%%%%%%%%%%%%%%%%%%%%%%%%%%%%%%%%%%%%%%%%%%%%%%%%%
\subsection*{Temperature Measurement Challenges}
\addcontentsline{toc}{subsection}{Temperature Measurement Challenges}
%%%%%%%%%%%%%%%%%%%%%%%%%%%%%%%%%%%%%%%%%%%%%%%%%%%%

\subsubsection*{Measurement Technologies}
\addcontentsline{toc}{subsubsection}{Measurement Technologies}

Determining accurate air temperature around components is complex due to:
\begin{itemize}
    \item Various sensor types with different characteristics:
    \begin{itemize}
        \item Thermocouples (T-type, E-type) using Seebeck effect with voltage output
        \item Resistance Temperature Detectors (RTDs) with excellent accuracy but requiring power
        \item Infrared sensors for surface temperature measurement
    \end{itemize}
    \item Calibration requirements for all sensors
    \item Reaction time variations based on sensor diameter and environmental conditions
    \item External factors affecting readings:
    \begin{itemize}
        \item Heat flux on the sensor
        \item Radiation losses/gains
        \item Conduction and convection effects
    \end{itemize}
\end{itemize}

\subsubsection*{Sensor Response Analysis}
\addcontentsline{toc}{subsubsection}{Sensor Response Analysis}

For accurate transient temperature measurements:
\begin{itemize}
    \item Sensor response time must be analyzed
    \item For thermal masses with Biot number $\ll$ 1, lumped capacitance model applies:
    \begin{itemize}
        \item Sensor behavior modeled as a single thermal mass
        \item Response characteristics calculated
    \end{itemize}
    \item Sensor diameter critically affects response time:
    \begin{itemize}
        \item Larger diameter sensors miss rapid temperature fluctuations
        \item 1mm probe significantly lags behind actual temperature changes
        \item Proper sensor selection vital for capturing transient thermal behavior
    \end{itemize}
\end{itemize}

%%%%%%%%%%%%%%%%%%%%%%%%%%%%%%%%%%%%%%%%%%%%%%%%%%%%
\subsection*{Convection Analysis}
\addcontentsline{toc}{subsection}{Convection Analysis}
%%%%%%%%%%%%%%%%%%%%%%%%%%%%%%%%%%%%%%%%%%%%%%%%%%%%

\subsubsection*{Forced Convection Assessment}
\addcontentsline{toc}{subsubsection}{Forced Convection Assessment}

To understand the convective boundary conditions:
\begin{itemize}
    \item Reynolds number analysis to determine flow regime:
    \[
    Re = \frac{\text{Inertial forces}}{\text{Viscous forces}} = \frac{\rho u L}{\mu}
    \]
    \item Nusselt number to relate convective to conductive heat transfer:
    \[
    Nu = \frac{\text{Total heat transfer}}{\text{Conductive heat transfer}} = \frac{h L}{k}
    \]
    \item Heat transfer coefficient can be determined from:
    \[
    h = \frac{Nu \cdot k}{L}
    \]
\end{itemize}

\subsubsection*{Flow Regime Identification}
\addcontentsline{toc}{subsubsection}{Flow Regime Identification}

Flow characterization guided the analysis approach:
\begin{itemize}
    \item Laminar flow (Re < 2300 in pipes)
    \item Transition/mixed convection (2300 < Re < 10,000)
    \item Fully turbulent flow (Re > 10,000)
    \item Factors affecting Nusselt number:
    \begin{itemize}
        \item Fluid properties
        \item Boundary layer development
        \item Surface roughness (Moody diagram)
    \end{itemize}
\end{itemize}

\subsubsection*{CFD vs. Correlations}
\addcontentsline{toc}{subsubsection}{CFD vs. Correlations}

Decision process for analysis method:
\begin{itemize}
    \item Computational Fluid Dynamics (CFD):
    \begin{itemize}
        \item Provides detailed fluid flow and heat transfer information
        \item Requires significant validation
        \item Demands substantial computational resources
        \item May be excessive for simpler problems
    \end{itemize}
    \item Standard correlations:
    \begin{itemize}
        \item More efficient if solution is relatively insensitive to exact heat transfer coefficient
        \item Useful when design is robust across a range of convection conditions
        \item For turbulent flow, Nusselt correlation can be used
    \end{itemize}
\end{itemize}

For the pressure sensor case, estimated heat transfer coefficients ranged from 17 to 40 W/m²K.

%%%%%%%%%%%%%%%%%%%%%%%%%%%%%%%%%%%%%%%%%%%%%%%%%%%%
\subsection*{Radiation Heat Transfer Analysis}
\addcontentsline{toc}{subsection}{Radiation Heat Transfer Analysis}
%%%%%%%%%%%%%%%%%%%%%%%%%%%%%%%%%%%%%%%%%%%%%%%%%%%%

\subsubsection*{Radiation Considerations}
\addcontentsline{toc}{subsubsection}{Radiation Considerations}

Radiation effects required evaluation due to:
\begin{itemize}
    \item Proximity to hot components (e.g., exhaust primaries)
    \item Heat rejection to environment
    \item Temperature-dependent effects (proportional to $T^4$)
    \item Emissivity variations across materials and temperatures
\end{itemize}

Key radiation heat transfer principles:
\begin{itemize}
    \item Stefan-Boltzmann law for radiation heat transfer
    \item Dependence on temperature to the fourth power
    \item View factor calculations between surfaces
    \item Advanced methods like Monte Carlo ray tracing when necessary
\end{itemize}

\subsubsection*{Material Emissivity Characterization}
\addcontentsline{toc}{subsubsection}{Material Emissivity Characterization}

Emissivity measurement is critical and complex:
\begin{itemize}
    \item Parameters affecting emissivity:
    \begin{itemize}
        \item Surface finish
        \item Material composition
        \item Temperature
        \item Viewing angle
        \item Wavelength of radiation
    \end{itemize}
    \item Specialized equipment for measurement:
    \begin{itemize}
        \item FTIR (Fourier Transform Infrared) spectroscopy
        \item Emissometers (handheld and laboratory grade)
        \item Colorimeters for total hemispheric emissivity
    \end{itemize}
    \item Characterization across spectral ranges and temperatures
    \item Materials with specialized emissivity properties:
    \begin{itemize}
        \item Carbon composites (low reflectance)
        \item Aerogels (insulation materials)
        \item Anodized aluminum
    \end{itemize}
\end{itemize}

%%%%%%%%%%%%%%%%%%%%%%%%%%%%%%%%%%%%%%%%%%%%%%%%%%%%
\subsection*{Component Temperature Distribution Analysis}
\addcontentsline{toc}{subsection}{Component Temperature Distribution Analysis}
%%%%%%%%%%%%%%%%%%%%%%%%%%%%%%%%%%%%%%%%%%%%%%%%%%%%

\subsubsection*{Thermal Imaging Assessment}
\addcontentsline{toc}{subsubsection}{Thermal Imaging Assessment}

Infrared thermography was used to:
\begin{itemize}
    \item Identify internal heating patterns within the sensor
    \item Determine temperature distribution uniformity
    \item Guide thermal management strategy (uniform insulation vs. targeted cooling)
\end{itemize}

Results showed:
\begin{itemize}
    \item Aluminum casing provided even heat distribution
    \item Small temperature variation (20.85°C to 29.3°C)
    \item No significant hot spots requiring targeted cooling
    \item In-plane temperature distribution more important than through-thickness behavior
\end{itemize}

\subsubsection*{Infrared Thermography Techniques}
\addcontentsline{toc}{subsubsection}{Infrared Thermography Techniques}

Proper thermography application required attention to:
\begin{itemize}
    \item Technology selection:
    \begin{itemize}
        \item Thermal cameras
        \item Infrared single-point sensors
    \end{itemize}
    \item Wavelength ranges:
    \begin{itemize}
        \item Short wave (0.7-3 $\mu$m)
        \item Medium wave (3-6 $\mu$m)
        \item Long wave (8-14 $\mu$m)
    \end{itemize}
    \item Measurement challenges:
    \begin{itemize}
        \item Surface coating limitations (to preserve thermal performance)
        \item Viewing through windows or hot gases
        \item Material-specific emissivity variations across wavelengths
    \end{itemize}
    \item Window materials for specific applications:
    \begin{itemize}
        \item Calcium fluoride
        \item Potassium bromide
        \item Selection based on transmission in target wavelength range
    \end{itemize}
\end{itemize}

%%%%%%%%%%%%%%%%%%%%%%%%%%%%%%%%%%%%%%%%%%%%%%%%%%%%
\subsection*{Thermal Conductivity Measurement}
\addcontentsline{toc}{subsection}{Thermal Conductivity Measurement}
%%%%%%%%%%%%%%%%%%%%%%%%%%%%%%%%%%%%%%%%%%%%%%%%%%%%

\subsubsection*{Custom Testing Apparatus}
\addcontentsline{toc}{subsubsection}{Custom Testing Apparatus}

For non-standard components like the pressure sensor:
\begin{itemize}
    \item Custom testing rigs were developed
    \item Hooke's law approach for thermal conductivity measurement:
    \begin{itemize}
        \item Heating element applied controlled heat flux
        \item Specimen placed between hot and cold plates
        \item Hot guards ensured one-dimensional heat flow
        \item Temperature differential measured across specimen
    \end{itemize}
    \item Thermal conductivity calculated based on heat flux and temperature gradient:
    \[
    k = \frac{q'' \cdot \Delta x}{\Delta T}
    \]
\end{itemize}

Testing revealed the sensor had:
\begin{itemize}
    \item Direction-dependent thermal conductivity
    \item Approximately 5 W/mK thermal conductivity in the measured direction
\end{itemize}

\subsubsection*{Thermal Interface Resistance}
\addcontentsline{toc}{subsubsection}{Thermal Interface Resistance}

Thermal contact resistance analysis:
\begin{itemize}
    \item Gaps between sensor and composite create thermal resistance
    \item Temperature drop occurs at interfaces
    \item Methods to minimize thermal contact resistance:
    \begin{itemize}
        \item Thermal interface materials
        \item Liquid metals
        \item Nano-engineered interface surfaces
    \end{itemize}
    \item Analysis determined whether thermal interface material was necessary
\end{itemize}

%%%%%%%%%%%%%%%%%%%%%%%%%%%%%%%%%%%%%%%%%%%%%%%%%%%%
\subsection*{Material Thermal Property Characterization}
\addcontentsline{toc}{subsection}{Material Thermal Property Characterization}
%%%%%%%%%%%%%%%%%%%%%%%%%%%%%%%%%%%%%%%%%%%%%%%%%%%%

\subsubsection*{Composite Materials Analysis}
\addcontentsline{toc}{subsubsection}{Composite Materials Analysis}

Carbon fiber reinforced polymers require special consideration:
\begin{itemize}
    \item Orthotropic thermal properties:
    \begin{itemize}
        \item Different thermal conductivity in-plane vs. through-thickness
        \item Potential anisotropy within the plane (x vs. y direction)
    \end{itemize}
    \item Measurement techniques:
    \begin{itemize}
        \item Laser Flash Analyzer (LFA) for thermal diffusivity
        \item Differential Scanning Calorimeter (DSC) for heat capacity
        \item Combined to calculate thermal conductivity:
        \[
        k = \alpha \cdot \rho \cdot c_p
        \]
        where $\alpha$ is thermal diffusivity, $\rho$ is density, and $c_p$ is specific heat
    \end{itemize}
\end{itemize}

Typical thermal conductivity values:
\begin{itemize}
    \item In-plane (x-y): approximately 3 W/mK
    \item Through-thickness (z): approximately 0.5 W/mK
\end{itemize}

\subsubsection*{Thermal Conductivity Range of Materials}
\addcontentsline{toc}{subsubsection}{Thermal Conductivity Range of Materials}

Materials used in F1 cover a wide spectrum of thermal conductivities:
\begin{itemize}
    \item Insulators:
    \begin{itemize}
        \item Multi-layer insulation (space exploration): near-zero conductivity
        \item Aerogels: 0.015 W/mK
        \item Home insulation (fiberglass): 0.03 W/mK
        \item Polymeric foams: 0.02-0.04 W/mK
        \item High-temperature insulators: 0.04-0.2 W/mK
    \end{itemize}
    \item Conductors:
    \begin{itemize}
        \item Polymers: 0.25 W/mK
        \item Copper: 400 W/mK
        \item Advanced materials (diamond, graphene, pyrolytic graphite): extremely high conductivity
    \end{itemize}
\end{itemize}

%%%%%%%%%%%%%%%%%%%%%%%%%%%%%%%%%%%%%%%%%%%%%%%%%%%%
\subsection*{Solution Implementation}
\addcontentsline{toc}{subsection}{Solution Implementation}
%%%%%%%%%%%%%%%%%%%%%%%%%%%%%%%%%%%%%%%%%%%%%%%%%%%%

\subsubsection*{Thermal Resistance Network Approach}
\addcontentsline{toc}{subsubsection}{Thermal Resistance Network Approach}

For the pressure sensor thermal management:
\begin{itemize}
    \item Thermal resistance network model developed as an alternative to full CFD
    \item Model incorporated all measured parameters:
    \begin{itemize}
        \item Convection coefficients
        \item Radiation properties
        \item Material thermal conductivities
        \item Interface resistances
    \end{itemize}
    \item Solution methods:
    \begin{itemize}
        \item Hand calculations
        \item Excel spreadsheets
        \item MATLAB and Simscape simulation
    \end{itemize}
    \item Model validated against car data
\end{itemize}

\subsubsection*{Final Solution and Results}
\addcontentsline{toc}{subsubsection}{Final Solution and Results}

After testing multiple configurations:
\begin{itemize}
    \item Encapsulated aerogel insulation implemented
    \item Temperature increase reduced from 31°C to 7°C above ambient
    \item Further improvements targeting 1°C above ambient:
    \begin{itemize}
        \item Space constraints present challenges
        \item Heat spreaders and heat pipes being considered
    \end{itemize}
\end{itemize}

This case study represents one of 20-30 similar thermal management tasks addressed annually.

%%%%%%%%%%%%%%%%%%%%%%%%%%%%%%%%%%%%%%%%%%%%%%%%%%%%
\subsection*{Brake System Thermal Management}
\addcontentsline{toc}{subsection}{Brake System Thermal Management}
%%%%%%%%%%%%%%%%%%%%%%%%%%%%%%%%%%%%%%%%%%%%%%%%%%%%

\subsubsection*{Evolution of Brake Design}
\addcontentsline{toc}{subsubsection}{Evolution of Brake Design}

Formula One brake designs have evolved substantially:
\begin{itemize}
    \item Earlier designs (2002): approximately 72 cooling drillings
    \item Current designs (2023): over 1,000 drillings of 2.5-3mm diameter
    \item Changes driven by:
    \begin{itemize}
        \item Increased vehicle speeds
        \item More demanding braking requirements
        \item Higher thermal loads
    \end{itemize}
\end{itemize}

\subsubsection*{Brake System Thermal Characteristics}
\addcontentsline{toc}{subsubsection}{Brake System Thermal Characteristics}

Extreme thermal conditions in F1 braking systems:
\begin{itemize}
    \item Heat loads up to 5 MW/m² during peak braking
    \item Carbon-carbon composite materials similar to aerospace applications
    \item Oxidation-limited materials:
    \begin{itemize}
        \item Higher temperatures increase wear rate
        \item Carbon converts to $\mathrm{CO}$ and $\mathrm{CO}_2$ at high temperatures
        \item Requires careful temperature management
    \end{itemize}
    \item Track-specific cooling requirements:
    \begin{itemize}
        \item Monaco vs. Austin vs. Indianapolis require different solutions
        \item Cooling must be optimized for specific energy profiles
    \end{itemize}
    \item Airflow minimization critical for aerodynamic performance
    \item Highly transient thermal behavior during race conditions
\end{itemize}

\subsubsection*{Carbon-Carbon Material Characterization}
\addcontentsline{toc}{subsubsection}{Carbon-Carbon Material Characterization}

Understanding brake disc material behavior:
\begin{itemize}
    \item Oxidation dependent on:
    \begin{itemize}
        \item Temperature
        \item Oxygen availability
        \item Gas flow velocity over surfaces
        \item Diffusion into porous material
    \end{itemize}
    \item Characterization techniques:
    \begin{itemize}
        \item Thermogravimetric Analysis (TGA)
        \item Thermomechanical Analysis (TMA)
        \item CT scanning for porosity analysis
        \item Custom test rigs for race-specific conditions
    \end{itemize}
    \item Material porosity:
    \begin{itemize}
        \item Typically 5-20\% porosity
        \item Increases with oxidation
        \item Affects structural integrity and thermal performance
        \item Varies through thickness and between discs
    \end{itemize}
\end{itemize}

\subsubsection*{Brake Thermal Properties Analysis}
\addcontentsline{toc}{subsubsection}{Brake Thermal Properties Analysis}

Full thermal characterization requires:
\begin{itemize}
    \item Density and porosity determination (CT scanning)
    \item Heat capacity measurement (DSC - independent of porosity)
    \item Thermal diffusivity testing:
    \begin{itemize}
        \item Laser Flash Analyzer (LFA)
        \item Challenges with porous materials and thickness measurement
        \item Location-dependent properties within the same disc
    \end{itemize}
    \item Emissivity characterization across temperature and wavelength ranges
    \item Combined properties for thermal simulations:
    \begin{itemize}
        \item 1D thermal transient simulations
        \item 3D steady-state CFD
        \item Oxidation and wear simulations
    \end{itemize}
\end{itemize}

\subsubsection*{Validation and Implementation}
\addcontentsline{toc}{subsubsection}{Validation and Implementation}

Comprehensive testing and simulation:
\begin{itemize}
    \item Dynamometer testing to validate models
    \item Temperature measurements at multiple disc locations
    \item Oxidation rate correlation with temperature and flow speed
    \item Pre-event simulations to determine optimal cooling configuration
    \item Wear prediction based on temperature profiles and oxidation rates
    \item Temperature limits established (e.g., avoiding 1100°C at high speeds to prevent excessive wear)
\end{itemize}

%%%%%%%%%%%%%%%%%%%%%%%%%%%%%%%%%%%%%%%%%%%%%%%%%%%%
\subsection*{General Insights on Thermal Engineering}
\addcontentsline{toc}{subsection}{General Insights on Thermal Engineering}
%%%%%%%%%%%%%%%%%%%%%%%%%%%%%%%%%%%%%%%%%%%%%%%%%%%%

\subsubsection*{Measurement Challenges}
\addcontentsline{toc}{subsubsection}{Measurement Challenges}

Key insights from Formula One thermal engineering:
\begin{itemize}
    \item Measurements are inherently difficult in thermal engineering
    \item Engineers often search for new physics to explain incorrect measurements
    \item Important to verify measurements before developing complex theories
    \item CFD provides answers but requires careful validation against reality
    \item Return to fundamentals (Reynolds number, heat transfer correlations) when results seem unreasonable
\end{itemize}

\subsubsection*{Thermal Engineering Career Perspective}
\addcontentsline{toc}{subsubsection}{Thermal Engineering Career Perspective}

Thermal engineering as a career field:
\begin{itemize}
    \item Offers diverse and exciting problem-solving opportunities
    \item Required across virtually all engineering industries
    \item Provides career flexibility and mobility
    \item Involves continuous learning and adaptation
    \item Presents new challenges regularly, preventing monotony
\end{itemize}

%%%%%%%%%%%%%%%%%%%%%%%%%%%%%%%%%%%%%%%%%%%%%%%%%%%%
\subsection*{Racing Dynamics Insights}
\addcontentsline{toc}{subsection}{Racing Dynamics Insights}
%%%%%%%%%%%%%%%%%%%%%%%%%%%%%%%%%%%%%%%%%%%%%%%%%%%%

\subsubsection*{Vehicle Performance Factors}
\addcontentsline{toc}{subsubsection}{Vehicle Performance Factors}

Relationship between car handling and driver performance:
\begin{itemize}
    \item Fast cars typically operate at the edge of stability
    \item Slower cars are generally easier to drive
    \item Superior drivers can handle challenging car dynamics
    \item Less experienced drivers may overdrive difficult cars:
    \begin{itemize}
        \item Braking too late
        \item Overheating tires
        \item Locking brakes
        \item Forcing turns beyond capability
    \end{itemize}
    \item Driver adaptability critical to extracting maximum performance
\end{itemize}

\subsubsection*{Aerodynamics and Thermal Integration}
\addcontentsline{toc}{subsubsection}{Aerodynamics and Thermal Integration}

Cross-functional collaboration in Formula One:
\begin{itemize}
    \item Thermal engineers provide cooling requirements
    \item Specific requirements for different components:
    \begin{itemize}
        \item Mass flow rates for clutch cooling
        \item Electronic component cooling needs
        \item Target temperatures for brake systems
    \end{itemize}
    \item CFD specialists run simulations to achieve required cooling
    \item Compromise between thermal and aerodynamic performance
    \item Design must account for "traffic margin":
    \begin{itemize}
        \item Cars rarely run in completely clean air
        \item Following other cars impacts cooling efficiency
        \item Managed through design or driver technique (temporary lift-off)
    \end{itemize}
    \item Track-specific optimizations performed for each race weekend
\end{itemize}